\section{Augmented inverses and the analytic family}
\label{sec:family}

We now reduce the metrics in Theorem~\ref{thm:tree-inequality} to a finite-dimensional family. The reduction keeps every kernel and cokernel direction as finite data. Uniform estimates across long seams will give a comparison error quadratic in the Ricci size.

\subsection{Weighted operators, finite augmentations, and seams}
\label{ell:inverses}

The weighted Fredholm framework on cylindrical ends, including
exceptional weights and separate weights at different ends, was
established by Lockhart and McOwen
\cite[Theorems~1.1, 6.2, and~8.1]{LockhartMcOwen1985}.
The finite-dimensional gluing reductions of Biquard and Ozuch are
close precedents for keeping the obstruction variables
\cite[arXiv version, Proposition~8.1]{Biquard2011}
\cite[arXiv version, Theorem~4.6]{Ozuch2022II}.
We prove the disconnected and seam estimates below in the precise
weights and finite trace spaces required for the later complex
continuation.

Choose
\[
 \begin{gathered}
 q_0=1+\tfrac12\bigl(\min(q,2)-1\bigr),\qquad
 q_1=\tfrac12(1+q_0),\\
 \lambda=1/n<\tfrac1{100}\min(\mu,1),\qquad
 \nu=\lambda/4,\quad \nu'=\lambda/2,\quad \gamma_g=2\lambda.
 \end{gathered}
\]
Here $n$ is an integer chosen once.  Decrease $\mu$ first if
necessary to account for any previous weakening of the end estimates.
We use metric Sobolev order $k=8$, orders $6$ through $9$ when taking
traces or variations, and one additional vector order for the gauge.
There is ample room in the assumed derivative bounds for these choices.

On an end use the cylindrical measure $d\tau\,d\theta$ and derivatives
$\partial_\tau$ and angular derivatives.  For metric tensors let
$X^k_\gamma$ be the sum of the ordinary core $H^k$ norms and the
$H^k$ norms of $e^{\gamma\tau}U$ on joined ends, together with the
root norm of $r^{q_0}U$.  Components are Cartesian components in
vertex units.  For a displacement vector use $V/r$ in the same
definition.  A second-order source space $Y^{k-2}_\gamma$ has two
fewer derivatives and includes multiplication of the equation by
the square of the local length scale; for vectors use also the
normalization by $r$.  These conventions make all coefficients
bounded on fixed-size cylindrical coordinate patches.  Finite
dimensional unknowns and tests carry fixed Euclidean norms.

Replace $h$ by a smooth background $b=b_H$ which agrees with $h$
up to a large fixed $\tau=H$ and is exactly flat for $\tau\ge H+1$;
make the corresponding fixed truncation of the root end.  The
background has no length dependence in vertex coordinates.  Its
defect tends to zero in the indicated spaces as $H\to\infty$.

The vector operator is the linearized tension of the identity map
with both metrics equal to $b$.  The metric operator is the
linearization of
\begin{equation}
 P_b(g)=\Ric(g)-\tfrac12\mathcal L_{W_b(g)}g,
 \qquad W_b(g)^i=g^{jk}\bigl(\Gamma(g)^i_{jk}-\Gamma(b)^i_{jk}\bigr).
 \label{ell:deturck}
\end{equation}
The principal part of $P_b$ is
$-\frac12g^{ab}\partial_a\partial_bg_{ij}$: the two derivatives of
the divergence of $g$ cancel.  On a flat end both linear operators
are component Laplacians, up to a fixed nonzero factor.

\begin{lemma}[Disconnected augmented inverses]
\label{ell:disconnected-inverse}
There are finitely many fixed smooth core-supported sources and
finitely many fixed linear tests on core compact sets such that
each disconnected operator, augmented by these sources and tests,
is an isomorphism $X^k_\gamma\oplus\R^a\to
Y^{k-2}_\gamma\oplus\R^b$.  The same choices work for the selected
orders and nearby sufficiently small positive joining weights.
They work for $b_H$, with uniformly bounded inverse, when $H$ is
sufficiently large.
\end{lemma}

\begin{proof}
For a spherical harmonic of degree $\ell$, a harmonic Cartesian
component has radial powers $r^\ell$ and $r^{-\ell-2}$.
Thus the exponents in $e^{\kappa\tau}$ are
\begin{center}
\begin{tabular}{lll}
\toprule
Components & Punctured end & Exterior end\\
\midrule
Metric & $-\ell,\ \ell+2$ & $\ell,\ -\ell-2$\\
Normalized vector $V/r$ & $1-\ell,\ \ell+3$ & $\ell-1,\ -\ell-3$\\
\bottomrule
\end{tabular}
\end{center}
Restrict these lists to the equivariant degree spaces.  The chosen
weights avoid every exponent, on the joined ends and on the root.
After factoring a harmonic channel into two first-order factors,
integrate each factor toward the end where its weighted exponential
decays.  Partial fractions supply a divisor
$|\kappa_1-\kappa_2|=2(\ell+1)$.  Except in the finitely many lowest
channels, the weighted kernels are bounded by
\begin{equation}
 \frac{C}{1+\ell}\exp\{-c(1+\ell)|\tau-\tau'|\}.
 \label{ell:kernel-bound}
\end{equation}
The finitely many other channels have a positive gap fixed by the
weight.  Integration and the equation give two derivatives of gain:
the kernel $L^1$ norm is $O((1+\ell)^{-2})$, and differentiation
uses one power of $1+\ell$ per derivative.  Summing squared
harmonic coefficients proves the Sobolev estimate.

Cut off these full-cylinder inverses on the ends and combine them
with local interior elliptic parametrices on the compact cores.
Commutators have one derivative less than the operator; after
restriction to a compact set their errors are compact by the
Rellich theorem.  The difference between the actual and flat
operators has arbitrarily small norm sufficiently far along each
end.  This gives both a left and a right parametrix modulo compact
operators and a small operator, hence a Fredholm operator.
Choose core tests separating its finite-dimensional kernel.
Compactly supported sources are dense in the source space, so
finitely many such sources span its finite-dimensional cokernel;
enlarge the designated core to contain their supports.  Splitting
off kernel and cokernel shows directly that adding these sources
as unknowns and adjoining the tests gives a bijection.  The bounded
inverse theorem gives the estimate.

The kernels and cokernels are unchanged when a small positive
weight moves without crossing a root.  On flat half-ends this
follows by separating modes; for the original decaying coefficient
perturbation, the same argument bootstraps decay up to the next
root.  This permits common choices at the finitely many weights.
Elliptic regularity makes the inverses agree at the different
Sobolev orders.  Finally, the operator difference for $b_H-h$
tends to zero: the same weight is present on input and output,
and all scale-normalized coefficient differences and their needed
derivatives tend to zero.  If $C_0$ bounds the chosen inverses,
choose this difference below $(2C_0)^{-1}$ and apply the Neumann
series.  The new inverse norm is at most $2C_0$.
\end{proof}

At a finite seam match the value and the opposite normal
$\tau$-derivatives of the normalized components.  Also specify
the common value in the critical spherical channel: degree zero
for metrics and degree one for vectors.  In problems with jumps
use the parent-side value for this last test.  For a vector,
$V/r=V_{\rm phys}/\rho$ in both charts, so these conditions match
the physical vector and its first derivatives.  The opposite
normal signs are forced by $\tau_p+\tau_c=2L_e$.

\begin{lemma}[Uniform inverse with finite seams]
\label{ell:seam-inverse}
Adjoin these seam data to the disconnected augmentations.  For all
sufficiently large real lengths the resulting piecewise problem
is an isomorphism, uniformly in the lengths.  Its data norms at
an edge are $e^{\gamma L_e}H^{k-1/2}$ for value jumps and critical
tests and $e^{\gamma L_e}H^{k-3/2}$ for derivative jumps.
The assertion includes any partial disconnection of the tree,
with the tests and matching omitted on opened edges.
\end{lemma}

\begin{proof}
Extend the source from each finite half boundedly across its
endpoint to the corresponding infinite end, and apply
Lemma~\ref{ell:disconnected-inverse}.  We correct the resulting
seam errors by all nondecaying modes, not just strictly growing
ones.  In a noncritical metric channel the incoming roots are
$\alpha=\ell+2$ on the point side and $\beta=\ell$ on the exterior
side.  For vectors they are $\alpha=\ell+3$ and $\beta=\ell-1$.
The vector degree-zero channel is absent: an invariant constant
vector would give a fixed point on $S^3$ for every nonidentity
element of a nontrivial freely acting group.

Normalize incoming amplitudes at the seam.  The leading map to
value and derivative jumps is
\[
 \binom{A-B}{\alpha A+\beta B},\qquad
 \begin{pmatrix}1&-1\\\alpha&\beta\end{pmatrix},
 \qquad \det=2(\ell+1).
\]
This inverse has exactly the orders required by the two trace
spaces.  In a critical channel there are a point-side constant,
a point-side growing mode, and an exterior-side constant.  For
their amplitudes $(A,B,C)$ the three data are
\begin{equation}
 (A+B-C,\ \kappa B,\ A+B),\qquad
 \begin{pmatrix}1&1&-1\\0&\kappa&0\\1&1&0\end{pmatrix},
 \quad \kappa=2\text{ or }4.
 \label{ell:critical-matrix}
\end{equation}
Its determinant is $\kappa\ne0$.

Cut each incoming wave $Ae^{\kappa(\tau-L_e)}$ off before the
fixed core.  Correct its equation error there with the disconnected
inverse.  This makes a globally homogeneous augmented-equation
variation, with its accompanying finite source coefficients.
A unit weighted seam amplitude has core error
$O(e^{-\gamma L_e})$ for a constant wave and still smaller error
for a growing wave.  On another far end the correction is a
decaying homogeneous sum, whose first allowed decay has a strict
gap beyond $\gamma$.  Thus its weighted seam traces tend to zero,
including when the receiving edge has a different length.  The
factor $e^{-(1+\ell)(L_e-H)}$ controls all higher harmonics.
The full trace map is therefore a vanishing-norm perturbation of
the preceding block maps, and is invertible.

For completeness, a homogeneous solution has on each flat half
a sum of incoming and outgoing modes.  Subtract the corrected
incoming waves with its incoming coefficients.  The remainder
continues with decay on the disconnected ends and has zero core
tests, so the disconnected inverse makes it zero.  The trace
orders of these coefficients follow from Cauchy data and the
root gap.  This proves injectivity as well as the stated bound.
The construction never uses an edge that has been opened, proving
the partial-disconnection assertion.  Piecewise $H^k$ is sufficient;
no higher normal matching is imposed in advance.
\end{proof}

\subsection{Complex lengths and a small analytic family}
\label{ell:complex-family}

\begin{lemma}[Contours and complex inverses]
\label{ell:complex-inverse}
For some $A,c_0>0$ the inverse of Lemma~\ref{ell:seam-inverse}
continues to
\begin{equation}
 \Re L_e>A,\qquad |\Im L_e|<c_0(\Re L_e)^2.
 \label{ell:quadratic-domain}
\end{equation}
The field norm is measured on increasing-real-part contours.
One strengthens the source norm on the portion
$u\ge\Re L_e/5$ and the seam norm by a fixed polynomial in
$1+\Re L_e$.  With these norms the inverse is uniform.
The same assertion holds for opened ends continued along such
contours and for every partial disconnection.
\end{lemma}

\begin{proof}
Write $\ell_e=\Re L_e$ and choose a fixed smooth function $\chi$
which is zero on $(-\infty,1/3]$ and one on $[9/10,\infty)$.
Use
\[
 \tau(u)=u+i(\Im L_e)\chi(u/\ell_e),\qquad u\le\ell_e.
\]
It is real before $\ell_e/3$, has endpoint $L_e$, and has
$\tau'=1$ near the endpoint.  Its derivatives satisfy
$|\tau'|\le C(1+\ell_e)$, $|\tau^{(j)}|\le C_j(1+\ell_e)^2$
for each of the fixed orders used here, and $|\tau'|\ge1$.
All bends lie in the exactly flat region after increasing $A$.

For a real root $\kappa$,
\[
 |e^{\kappa(\tau(u)-\tau(u'))}|
       =e^{\kappa(u-u')}.
\]
Thus Equation~\eqref{ell:kernel-bound} remains valid in modulus.
The integration element $d\tau'$ and finitely many $u$ derivatives
cost only a polynomial in $1+\ell_e$.  The root gap still gives
the two derivative gain, with local terms controlled by the
equation.  Fix the polynomial exponent by bounding each
differentiation, product, trace, and extension appearing in this
proof by its elementary polynomial degree, and choosing $P$ to
be one plus the sum of those finitely many nonnegative degrees.
Only the fixed Sobolev orders enter this choice; no asymptotic
expansion order enters it.  Choose a fixed smooth cutoff
$\chi_0$ equal to zero for $u\le1/6$ and to one for $u\ge1/5$,
and put $\chi_{\ell_e}(u)=\chi_0(u/\ell_e)$.  On this half use
the strengthened source norm
\[
 \|f\|_{Y^{k-2}_\gamma}
       +(1+\ell_e)^P\|\chi_{\ell_e}f\|_{Y^{k-2}_\gamma},
\]
and multiply the seam-data norms by $(1+\ell_e)^P$.
The first summand retains control through the transition.
Indeed, decompose $f=(1-\chi_{\ell_e})f+\chi_{\ell_e}f$.
The first source is supported where $u\le\ell_e/5$, and its
Sobolev norm is bounded by the basic source norm, since the
cutoff derivatives are uniformly bounded.  Every transfer
from its support to the bend, which starts at $\ell_e/3$,
gains $e^{-\delta_0(2\ell_e/15)}$ for a fixed weighted root
gap $\delta_0>0$; this pays any polynomial contour loss.
On the unbent portion there is no such loss.  The second
summand of the strengthened norm pays the polynomial losses
for the remaining source.

To use a disconnected core inverse, solve forced data in the flat
part by these kernels, cut that solution off on a real collar,
and apply the real disconnected inverse to the remaining compact
error.  Its decaying homogeneous far field is continued by the
same exponentials.  This also constructs the inverse on an opened
end with a bend followed by a return to the real ray farther out.
If a bend starts increasingly far out, a cutoff for the forced
solve can be moved with it inside the flat region.  The real
core estimate remains the fixed estimate, and the flat estimates
are unchanged.  This proves the needed uniformity for such bends.

Extend finite-half data through the seam and apply this
disconnected construction.  Correct traces by the incoming modes
from Lemma~\ref{ell:seam-inverse}.  They have polynomially bounded
norms relative to their weighted seam amplitudes.  Their leading
trace matrices are unchanged because $\tau'=1$ at the endpoint.
The core correction is exponentially small.  Its outgoing
homogeneous traces at every receiving seam have a strict extra
decay rate, absorbing the polynomial at that receiving length.
For the initially forced solution, data before $\ell_e/5$ reach
the seam with exponential gain; data after that point have the
strengthened norm on that very end.  Increase $P$ once to pay
the subsequent trace correction as well.  This proves a uniform
bound and a vanishing perturbation of the block trace map.
The same representation proves injectivity.

Finally, a coefficient perturbation with size $O(e^{-\sigma u})$
is small in the strengthened operator norm on a bent region,
since $(1+\ell_e)^Pe^{-\sigma\ell_e/5}\to0$.
Small perturbations on the remaining real portions are handled
by the real inverse.  A Neumann series proves the perturbative
version that will be used below.  This argument is unchanged if
any subset of the ends is left open.
\end{proof}

Here and below ``uniform'' permits increasing the fixed lower
length threshold.  The choices have the following order, which
avoids a circular truncation requirement.  Fix weights, derivative
orders, the augmented operators at $h$, and their inverse bound
$C_0$.  Fix a contraction radius $\delta>0$ so that the product
and inverse constants give a nonlinear Lipschitz constant at most
$1/4$ on the $2\delta$ ball.  Choose $H$ so that the operator
perturbation is below $(2C_0)^{-1}$ and the reference defect is
below $\delta/(8C_0)$.  Choose the finite parameter radius so
that its additional defect has the same bound.  Finally choose
$A>10(H+1)$ large enough that all trace perturbations have norm
at most $1/4$ and that the required polynomial times exponential
bounds hold.  The finite number of partial disconnections permits
the same choices for all of them.  The inverse constants used to
choose $\delta$ are uniform in $H$ by
Lemma~\ref{ell:disconnected-inverse}.  The scale-normalized tensors
$b_H$ have a uniform positive-definite lower bound, and $b_H$,
$b_H^{-1}$ and their derivatives through the fixed orders used
here have bounds independent of $H$: the transition uses a
fixed-width cutoff translated into regions where $h$ approaches
$\delta$.  Consequently the local product and inverse-matrix
constants used to choose $\delta$ are independent of $H$ as well.
No new smallness radius is
chosen after increasing the expansion order.

\begin{lemma}[Gauge for the real sequence]
\label{ell:gauge}
After a diffeomorphism of the invariance class in
Lemma~\ref{ell:functional}, each sufficiently late metric in the
sequence satisfies
\begin{equation}
 W_b(g)=m\cdot\psi,
 \label{ell:gauge-equation}
\end{equation}
where $\psi$ is the finite list of vector augmentation sources,
rescaled in physical units, and $m$ is small.  The gauged metric
retains Equation~\eqref{eq:tree-ricci-error} in weaker admissible
weights and is close to $b$ at the joining weight $\gamma_g$.
\end{lemma}

\begin{proof}
Solve for a map $f:(N,g)\to(N,b)$ near the identity with
$\tau(f;g,b)=-m\cdot\psi(f)$.  Parametrize it by the background
exponential map of a displacement $V$, and impose zero core and
critical-seam tests on $V$.  The derivative in $(V,m)$ is the
augmented vector isomorphism.  On every coordinate patch of scale
$r$, this is the ordinary tension equation with bounded smooth
target coefficients.  Sobolev multiplication at the chosen
orders, with locally comparable weights, bounds its quadratic
remainder and its Lipschitz difference.  The uniform inverse and
the smallness of $g-b$ give a contraction.  The latter smallness
follows from compact convergence and the integrable dominating
tail $e^{-(\mu-\gamma_g)\tau}$; the root argument uses $q-q_0>0$.

Smallness of $V/r$ in scaled $C^1$ makes the entire displacement
isotopy locally invertible.  It is proper, since the displacement
is small relative to radius and decays on the root, so it is a
degree-one covering and hence a diffeomorphism.  The equation in
the image metric reads $-W_b(f_*g)=-m\cdot\psi$.
Seam matching of $V/r$ and its first derivatives gives a genuine
solution across the seam; the smooth equation then supplies the
higher compatibility.  At infinity $V=O(r^{1-q_0})$, which is
admissible for Lemma~\ref{ell:functional}.  Small scaled coordinate
changes preserve the end estimates after decreasing exponents.
\end{proof}

We now write $g$ for the gauged metric.  Choose the metric
augmentation sources $\varphi$ from
Lemma~\ref{ell:disconnected-inverse} and solve
\begin{equation}
 P_b(h_*)+\tfrac12\mathcal L_{m\cdot\psi}h_*=d\cdot\varphi,
 \qquad h_*=b+U.
 \label{ell:family-equation}
\end{equation}
We make the scaling of the augmentation sources explicit. On a
vertex of physical length $a_v$, let $D_{a_v}(x)=a_vx$ denote
passage from vertex to physical coordinates. A vector source in the
tension equation has physical components $a_v^{-1}\psi_v$:
its pullback as a vector field is $a_v^{-2}\psi_v$.
A covariant two-tensor source has physical components
$a_v^{-2}\varphi_v$, so $D_{a_v}^*\varphi_{\mathrm{phys}}=\varphi_v$.
Consequently, writing $k_v$ for a metric in vertex coordinates,
\[
 D_{a_v}^*\bigl(\mathcal L_{m\psi_{\mathrm{phys}}}k_{\mathrm{phys}}\bigr)
       =\mathcal L_{m\psi_v}k_v
 \quad\text{when }D_{a_v}^*k_{\mathrm{phys}}=a_v^2k_v.
\]
These are equation-source scalings; a displacement itself scales as
$a_vV_v$. With these conventions the normalized bulk equations and
core tests have no length dependence.

Here $d$ is an unknown finite vector.  The core test values for
$U$, the vector $m$, and the coefficients of the critical seam
values
\begin{equation}
                  \pi U|_e=e^{-\lambda L_e}z_e
 \label{ell:seam-prescription}
\end{equation}
make up a finite parameter vector $z$.  Values and first
derivatives genuinely match on every remaining seam.  On an
opened edge omit its matching and its critical prescription.

\begin{lemma}[Analytic small branch]
\label{ell:analytic-family}
On a fixed complex polydisc in $z$ and the quadratic length domain
of Equation~\eqref{ell:quadratic-domain},
Equation~\eqref{ell:family-equation} has a unique small solution
$(U,d)$.  It obeys uniform bounds at both joining weights $\nu$
and $\nu'$.  Its augmented linearization has the same uniform
inverse bounds.  These assertions hold for every partial
disconnection on common smaller domains.  The family, its compact
observables, and its endpoint values are single-valued holomorphic
functions of the length parameters.
\end{lemma}

\begin{proof}
On cores, inverse metric matrices and the tensor operations in
Equation~\eqref{ell:family-equation} are analytic in $U$ near zero;
the fixed spatial coefficients need only be smooth.  On an
exactly flat end, in neutralized Cartesian components,
\begin{equation}
                   \mathcal P U=\mathcal N(U),
 \label{ell:flat-equation}
\end{equation}
where $\mathcal P$ is the flat cylindrical component Laplacian
and $\mathcal N$ vanishes to second order.  Its Taylor operations
contain at most two derivatives and have polynomial angular
coefficients.  Indeed $r\partial_{x_i}$ is a combination of
$\partial_\tau$ and tangential derivatives with polynomial
coefficients in the unit Cartesian coordinates.  Consequently
finite angular harmonic sums remain finite at each Taylor order.

Apply Lemma~\ref{ell:complex-inverse} to the equation and all
tests.  A quadratic term has one additional factor of decay
relative to the source weight.  Its basic source norm is
bounded by $C\delta^2$, and the additional strengthened
contribution is bounded by
\[
 C\delta^2(1+\ell_e)^P e^{-\nu\ell_e/6},
\]
because the cutoff support lies in $u\ge\ell_e/6$.
For differences the corresponding bound replaces $\delta^2$
by $\delta\|U-V\|_{X^k_\nu}$.  Taking the supremum over
$\ell_e$ gives constants independent of the truncation and
of all other lengths.
The seam data have size at most a polynomial times
$e^{-(\lambda-\nu)\ell_e}|z_e|$.  The order of choices stated
above makes the equation a uniform contraction, at both $\nu$
and $\nu'$.  Uniqueness identifies the solutions.  A Neumann
series around the augmented linear operator proves invertibility
of the branch linearization.  Analytic convergence of the
contraction proves analyticity in ordinary parameters and in
local holomorphic variations of endpoints on a fixed contour.

We specify why this is analyticity in $L$, despite a smooth rather
than holomorphic global choice of contours.  In a local
parametrization of a flat contour, a holomorphic endpoint change
changes the coefficients analytically through
$\partial_\tau=(d\tau/du)^{-1}\partial_u$.
For an infinitesimal parametrization change $v(u)=\delta\tau(u)$
fixing the endpoints, the candidate variation $v\partial_\tau U$
satisfies the differentiated equations.  The normal operator varies
by $\delta\partial_\tau=-(\partial_\tau v)\partial_\tau$, so
\[
 \delta(\partial_\tau U)
 =\partial_\tau(v\partial_\tau U)
   -(\partial_\tau v)\partial_\tau U
 =v\partial_\tau^2U.
\]
Thus $v=0$ at each endpoint preserves both matching conditions and
the critical-value tests; on the fixed core portions it also preserves
the core tests.  This variation belongs to $X^{k-1}_\nu=X^7_\nu$.
The lower-order inverse already supplied by
Lemmas~\ref{ell:disconnected-inverse}--\ref{ell:complex-inverse}
uses the same sources and tests, with value and critical traces in
$H^{13/2}$ and normal traces in $H^{11/2}$.  The same small-perturbation
argument gives this inverse for the branch linearization.  Its
uniqueness therefore identifies $v\partial_\tau U$ with the actual
solution variation, with zero variation of the compact unknowns.
Integrating it through a homotopy shows contour independence
where the coordinate is fixed.  Monotone-real-part contours with
fixed endpoints are deformable through the same class; this
identifies the local holomorphic continuations and excludes
monodromy.  For an open end take a bend returning to the real ray;
deformations are compactly supported in its flat region.  Tail
limits on the unchanged ray are consequently unchanged.  An
autonomous analytic density integrated with $d\tau$ changes by
the endpoint term $v$ times that density, and therefore its
integral is also invariant under fixed-endpoint deformations.
The same shape calculation applies to linearized solutions.
\end{proof}

\begin{lemma}[Comparison with the actual metric]
\label{ell:comparison}
Choose $z$ from the core tests and gauge coefficients of $g$, and
from Equation~\eqref{ell:seam-prescription} at its seams.  Then
\begin{equation}
 \|g-h_*\|_{X^k_\nu}\le CM,\qquad
 \|\Ric(h_*)\|_{Y^{k-2}_\nu}\le CM,\qquad
 |E(g)-E(h_*)|\le CM^2.
 \label{ell:comparison-estimates}
\end{equation}
\end{lemma}

\begin{proof}
The seam parameters are small because $g-b$ has weight
$\gamma_g>\lambda$.  All parameters lie strictly inside the
chosen small polydisc by taking the truncation and sequence
thresholds sufficiently large.  In the gauge
Equation~\eqref{ell:gauge-equation}, the actual metric solves
Equation~\eqref{ell:family-equation} with $d=0$ and the extra source
$\Ric(g)$.  Its source norm is $O(M)$ by
Equation~\eqref{eq:tree-ricci-error}.  Subtraction and the uniform
inverse, absorbing the small nonlinear Lipschitz term, prove the
first estimate, including the difference in compact unknowns.
The local Ricci map from $X^k_\nu$ to $Y^{k-2}_\nu$ is Lipschitz
on the small ball, proving the second estimate.

Along the segment from $g$ to $h_*$ the Ricci source and the
metric variation are both $O(M)$.  In the interior the pairing
in Equation~\eqref{ell:first-variation} costs integrals of
$\rho^2$ against products of the joining weights; these are
bounded by the scale calculation in Lemma~\ref{ell:functional}.
On the root the radial integral is
$\int r^{1-2q_0}\,dr<\infty$.  Integrating the first variation
proves the third estimate.
\end{proof}

Define real functions on the real family by
\begin{equation}
 F=E(h_*),\qquad
 J=\int|\Ric(h_*)|_{h_*}^2\,d\mu_{h_*}
 +\int_{\substack{\text{root end}\\r>r_0}}
       r^{2q_1}|\Ric(h_*)|_\delta^2\,dx.
 \label{ell:observables}
\end{equation}
For complex metrics the squares are analytic bilinear
contractions, without complex conjugation.  These give bounded
holomorphic continuations of $F,J$ on the quadratic domains.
For $F$, integrate the quadratic subtracted scalar density on a
fixed root tail and retain its linear boundary contribution at
$r_0$.  For $J$, the root weighted radial bound is
$\int r^{2q_1-2q_0-1}\,dr<\infty$.  Interior Ricci-square
integrals are scale invariant in dimension four and integrable
in cylindrical measure.  Scalar-volume integrals carry the factor
$a_v^2r^2$: on a child half this is
$a_p^2e^{-4L_e}e^{2\tau}$, whose modulus stays bounded.
Every contour polynomial is absorbed by the available exponential
decay.  On the approximants in Lemma~\ref{ell:comparison},
\begin{equation}
                  0\le J\le CM^2.
 \label{ell:J-small}
\end{equation}
